\input{Preamble}%

\begin{document}
\ctikzset{european, straight voltages, cute inductors, bipoles/length=1cm}%
\begin{tikzpicture}[%
    show background rectangle,%
    tight background,%
    background rectangle/.style={fill=white}%
  ]
  % Régulateur de tension triphasé en étoile : une unité par phase, de borne commune le neutre

  %
  % Définitions
  %
  % Une unité : autotransformateur à prise variable. Arguments : abscisse, ordonnée, étiquette.
  % L'enroulement va de la borne commune en (x,y-1.2) au point haut (x,y). L'entrée arrive à
  % gauche sur la prise en (x-0.8,y). L'admittance magnétisante Ym part de l'entrée en (x-1.8,y)
  % et rejoint le fil de la borne commune sous l'enroulement en (x,y-1.5). L'impédance série Z2
  % va du haut de l'enroulement à la sortie de l'unité en (x+1.6,y).
  \newcommand{\regulatorunit}[3]{%
    \draw[-{Latex[length=2.5mm]}] ({#1-0.8},#2) -- ({#1-0.12},{#2-0.5});%
    \node[below left] at ({#1-0.5},{#2-0.25}) {#3};%
    \draw ({#1},{#2-1.2}) to[cute inductor] ({#1},#2);%
    \draw ({#1-1.8},#2) to[generic, l_=$\underline{Y_{\mathrm{m}}}$, *-] ({#1-1.8},{#2-1.5})
    to[short, -*] ({#1},{#2-1.5});%
    \draw ({#1},#2) to[generic, l=$\underline{Z_2}$] ({#1+1.6},#2);%
  }%
  \pgfmathsetmacro{\r}{0.15};% Rayon des sauts de fils
  \pgfmathsetmacro{\xl}{0};%
  \pgfmathsetmacro{\xa}{3.5};%
  \pgfmathsetmacro{\xb}{7};%
  \pgfmathsetmacro{\xc}{10.5};%
  \pgfmathsetmacro{\xr}{13};%

  \pgfmathsetmacro{\yn}{0};%
  \pgfmathsetmacro{\yc}{2.5};%
  \pgfmathsetmacro{\yb}{5};%
  \pgfmathsetmacro{\ya}{7.5};%

  %
  % Dessin
  %
  % Tensions amont
  \node[left] at (\xl,\ya) {$\underline{V_{1,\arm}}$};
  \node[left] at (\xl,\yb) {$\underline{V_{1,\brm}}$};
  \node[left] at (\xl,\yc) {$\underline{V_{1,\crm}}$};
  \node[left] at (\xl,\yn) {$\underline{V_{1,\nrm}}$};

  % Tensions aval
  \node[right] at (\xr,\ya) {$\underline{V_{2,\arm}}$};
  \node[right] at (\xr,\yb) {$\underline{V_{2,\brm}}$};
  \node[right] at (\xr,\yc) {$\underline{V_{2,\crm}}$};
  \node[right] at (\xr,\yn) {$\underline{V_{2,\nrm}}$};

  % Unités
  \regulatorunit{\xa}{\ya}{$a_\anrm$}%
  \regulatorunit{\xb}{\yb}{$a_\bnrm$}%
  \regulatorunit{\xc}{\yc}{$a_\cnrm$}%

  % Entrées (les fils des phases b et c sautent par-dessus les bornes communes des autres unités)
  \draw (\xl,\ya) to[short,*-,i=$\underline{I_{1,\arm}}$] (\xl+1.2,\ya) -- (\xa-0.8,\ya);%
  \draw (\xl,\yb) to[short,*-,i=$\underline{I_{1,\brm}}$] (\xl+1.2,\yb)
    -- (\xa-\r,\yb) arc[start angle=180, end angle=0, radius=\r] -- (\xb-0.8,\yb);%
  \draw (\xl,\yc) to[short,*-,i=$\underline{I_{1,\crm}}$] (\xl+1.2,\yc)
    -- (\xa-\r,\yc) arc[start angle=180, end angle=0, radius=\r]
    -- (\xb-\r,\yc) arc[start angle=180, end angle=0, radius=\r] -- (\xc-0.8,\yc);%

  % Sorties
  \draw (\xa+1.6,\ya) -- (\xr-1.2,\ya) to[short,-*,i<=$\underline{I_{2,\arm}}$] (\xr,\ya);%
  \draw (\xb+1.6,\yb) -- (\xr-1.2,\yb) to[short,-*,i<=$\underline{I_{2,\brm}}$] (\xr,\yb);%
  \draw (\xc+1.6,\yc) -- (\xr-1.2,\yc) to[short,-*,i<=$\underline{I_{2,\crm}}$] (\xr,\yc);%

  % Neutre, borne commune des unités, transmis d'un côté à l'autre
  \draw (\xl,\yn) to[short,*-,i=$\underline{I_{1,\nrm}}$] (\xl+1.2,\yn)
    -- (\xr-1.2,\yn) to[short,-*,i<=$\underline{I_{2,\nrm}}$] (\xr,\yn);%
  \draw (\xa,\ya-1.2) to[short,-*] (\xa,\yn);%
  \draw (\xb,\yb-1.2) to[short,-*] (\xb,\yn);%
  \draw (\xc,\yc-1.2) to[short,-*] (\xc,\yn);%

\end{tikzpicture}
\end{document}
% Local Variables:
% mode: latex
% TeX-engine: luatex
% TeX-source-correlate-method-active: synctex
% ispell-local-dictionary: "british"
% coding: utf-8
% LaTeX-indent-level: 2
% fill-column: 100
% End:
